\begin{example}[title = Minimum independent set]
    e.g. We consider this task as a set of problems. \vspace{7pt}

    \textbf{Problem}. \\
    It asks whether a pair $(G, k)$ of an undirected graph and a natural number $k \in N$ is such that $G = (V, E)$ admits an independent set $W \subseteq V$ of 
    cardinality at least $k$. \vspace{7pt}

    \textbf{Note}. \\ 
    We are not computing the paramater $k$ but, instead, we provide it as an input parameter. We use it in order to check if it exists an independent set 
    of nodes of cardinality at least equal to $k$. \vspace{7pt}

    \textbf{Certificate}. \\ 
    In this problem or language, the certificate coincides with the independent set $W$. Check its cardinality and independence is quite easy, but what is really 
    tricky of this task is to find out the independent set $W$. \vspace{7pt}

    Formally, the certificate is described as follows:
    $$L = \{(G,k)| \exists W \ . W \ \text{is independent for} \ (G,k) \wedge |W| \geq k \wedge W \subseteq V\}$$
\end{example}
\begin{example}[title = Subset sum]
    e.g. Subset sum problem. \vspace{7pt}

    \textbf{Problem}. \\ 
    It asks whether, given a sequence of natural numbers $n_1, \dots, n_m$ and a number $k$, there exists a subset $I$ of \{1, \dots, m\} such that:
    $$ \sum_{i \in I} n_i = k $$ 

    \textbf{Certificate}. \\ 
    It's pretty simple to identify the certificate of this language: the certificate here is just $I$, checking whether the summation is equal to a natural number
    is very easy. \vspace{7pt}

    In this problem, checking if the summation of natural numbers is equal to $k$ is quite easy, but what is really tricky is to find out the subset of 
    natural numbers $I$.
\end{example}
\begin{example}[title = Composite Numbers]
    e.g. Composite numbers problem. \vspace{7pt}

    \textbf{Problem}. \\
    Given a number $n \in N$, determine whether $n$ is composite, not prime. \vspace{7pt}

    \textbf{Note}. \\
    In order to determine whether a number is composite or not, we don't need to take all its range of numbers, we can just use the $\sqrt{n}$. Anyway, 
    even though by this consideration we have reduced the range of numbers, we have to check any of them. \vspace{7pt}
    
    For sure, checking every single number up to $\sqrt{n}$ cannot be solved in polynomial time: the total amount of steps is exponential in the size of the input, which is $\sqrt{n}$. \vspace{7pt}

    \textbf{Certificate}. \\
    It is the factorization of $n$, a couple of natural numbers $(m,l)$ such that $n = m \cdot l$. \vspace{7pt}

    More in details, the verifier and the certificate are respectively:
    \begin{itemize}
        \renewcommand{\labelitemi}{-}
        \item Verifier, product between the two natural numbers, $n = m \cdot l$.
        \item Certificate, decomposition of $n$ into a pair of natural numbers.
    \end{itemize}
\end{example}
\begin{example}[title = Decisional Linear Programming]
    e.g. Decisional linear programming problem. \vspace{7pt}

    \textbf{Problem}. \\
    Given a sequence of $m$ linear inequalities with rational coefficients over $n$ variables, decide whether there is a rational assignment to the variables 
    $x_1, \dots, x_n$ which makes the inequalities true. \vspace{7pt}

    Inequalities are of the form:
    $$\sum_{i = 1}^{n} a_i x_i \leq b$$
    where:
    \begin{itemize}
        \renewcommand{\labelitemi}{-}
        \item $a_i \rightarrow$ rational coefficient
        \item $b_i \rightarrow$ rational number
    \end{itemize} \vspace{7pt}

    \textbf{Note}. \\
    As usual, in one hand checking the satisfiability of these inequalities is trivial. In the other hand, finding the $n$ different rational numbers for the 
    $x_1, \dots, x_n$ variables is difficult. \vspace{7pt}

    \textbf{Certificate}. \\
    The assignments, which can be easily checked for correctness. \vspace{7pt}

    Additionally, the assignments must be small enough, since the inequalities are linear. In the formal definition, this means that the binary encodings of the
    assignments taken in input by the Turing Machine cannot be too big.
\end{example}
\begin{example}[title = Decisional 0/1 Linear Programming]
    e.g. Decisional 0/1 linear programming problem. \vspace{7pt}

    \textbf{Problem}. \\
    Given a sequence of $m$ linear inequalities, decide whether there is an assignment of zeros and ones to the variables $x_1, \dots, x_n$ making the 
    inequalities true. \vspace{7pt}

    \textbf{Note}. \\
    We have just reduced the search space of the assignments, so the certificate still be the same as in the example above. However, making this problem 
    more "easy", we cannot anymore exploit some geometric properties of the inequalities to find out the best matches. \vspace{7pt}

    \textbf{Certificate}. \\
    Again, the assignments. 
\end{example}
Some of the discussed problems are also in \textbf{P}:
\begin{enumerate}
    \item Composite numbers.
    \item Decisional linear programming.
\end{enumerate}
All the other problems, as minimum independent set, subset sum and decisional 0/1 linear programming, are currently \textbf{not known} to be in P: this does 
not mean that they cannot be solved, but up to now does not exist any algorithm that can solve them in polynomial time. \vspace{7pt}